% Generated draft: requires mathematical and source review.
\documentclass[11pt]{article}
\usepackage[margin=25mm]{geometry}
\usepackage{fontspec}
\setmainfont{Latin Modern Roman}
\usepackage{amsmath,amssymb,mathtools,unicode-math}
\setmathfont{Latin Modern Math}
\usepackage{enumitem,needspace,xurl,hyperref}
\setlength{\parindent}{0pt}
\setlength{\parskip}{5pt}
\setlength{\emergencystretch}{3em}
\newcommand{\N}{\mathbb{N}}
\newcommand{\Z}{\mathbb{Z}}
\newcommand{\Q}{\mathbb{Q}}
\newcommand{\R}{\mathbb{R}}
\newcommand{\C}{\mathbb{C}}
\newcommand{\F}{\mathbb{F}}
\newcommand{\A}{\mathbb{A}}
\newcommand{\PP}{\mathbb{P}}
\newcommand{\E}{\mathbb{E}}
\newcommand{\eps}{\varepsilon}
\newcommand{\dd}{\,\mathrm{d}}
\begin{document}
\section*{UnsolvedMath ID: 30005504}
\textbf{Source problem code:} catalogue entry
\subsection{Definitions and mathematical statement}
Let \(G\) be a finite group. Write
\[
Z(G)=\{z\in G:zx=xz\text{ for every }x\in G\}
\]
for its center, and let \(\operatorname{Irr}(G)\) denote the irreducible complex characters of \(G\). For \(\chi\in\operatorname{Irr}(G)\), its degree is \(\chi(1)\); the character is nonlinear when \(\chi(1)>1\). A finite group \(H\) is quasisimple if \(H=[H,H]\) and \(H/Z(H)\) is nonabelian simple. The supplied record does not define its convention for \emph{almost quasisimple}; the intended class of groups must therefore be confirmed before this becomes a fully specified statement.

For every finite almost quasisimple group \(G\), in the source's intended sense, and every \(g\in G\setminus Z(G)\), the requested bounds are as follows.
\begin{enumerate}
\item[(A)] Find an explicit real number \(\gamma(g)\), as small as possible, such that
\[
0<\gamma(g)<1,
\qquad
\frac{|\chi(g)|}{\chi(1)}\leq\gamma(g)
\quad\text{for every }\chi\in\operatorname{Irr}(G)
\text{ with }\chi(1)>1.
\]
\item[(B)] Find an explicit real number \(\alpha(g)\), as small as possible, such that
\[
0<\alpha(g)<1,
\qquad
|\chi(g)|\leq\chi(1)^{\alpha(g)}
\quad\text{for every }\chi\in\operatorname{Irr}(G).
\]
\end{enumerate}
Here \(G\) is fixed when writing \(\gamma(g)\) and \(\alpha(g)\). Each bound must hold simultaneously for all characters in its stated range; the record does not specify which group or element invariants an explicit expression may use.

\emph{Unresolved qualifications.} The original excerpt asks for bounds that are ``as small as possible,'' whereas the cleaned statement asks for the ``smallest explicit'' bounds. No definition of explicitness, admissible family of expressions, or guarantee that a minimum exists within the open interval \((0,1)\) is supplied. The optimization request is retained without asserting existence of such minima. The precise convention for almost quasisimple groups, and whether the stated noncentrality condition ensures the requested strict bounds for every character in the indicated ranges, require mathematical review. No faithfulness hypothesis or restriction on character kernels is supplied.

\emph{Source context, not independently verified.} The excerpt reports that a combination of results attributed to Gluck, Gluck--Magaard, and Guralnick with the unnamed speaker permits \(\gamma=79/80\), unless \(E(G)=A_n\). Here \(E(G)\) denotes the layer, generated by the subnormal quasisimple subgroups of \(G\), and \(A_n\) denotes an alternating group. The precise hypotheses and meaning of the displayed exceptional case must be checked against the report; the supplied record does not give the referenced bibliography entries.

The excerpt then focuses on finite classical groups \(G=\operatorname{Cl}(V)\), with \(V=\F_q^n\), and discusses support using codimensions of kernels of \(g-\lambda I\) after scalar extension. The scalar-field notation and the numerical support-dependent character bound are damaged in the supplied transcription, so no precise auxiliary formula is asserted here. That auxiliary passage ranges over nontrivial irreducible characters and all elements of the classical group, a scope distinct from the two requested bounds above.
\subsection{Short English statement}
For each finite almost quasisimple group and each noncentral element, seek explicit optimal bounds on all nonlinear irreducible character ratios and on all irreducible character values as powers of their degrees. The definition of almost quasisimple, feasibility of strict bounds, and meaning and attainment of optimality require clarification. This is a statement-only draft requiring human mathematical review.
\subsection{Sources}
\begin{itemize}
\item Local UnsolvedMath registry, permanent ID \(30005504\), source code \texttt{catalogue entry}, title: \emph{Optimal Character-Ratio Bounds for Almost Quasisimple Groups}. The supplied record belongs to the Oberwolfach Reports Open Problems collection and identifies the excerpt as Problem 1. The project registries are \texttt{data/unsolvedmath/problems.json} and \texttt{catalog/problems.json}; this draft uses the supplied complete record without independently inspecting those files.
\item \emph{Representations of Finite Groups} (2023), Oberwolfach Reports, DOI 10.4171/owr/2023/19, pp.\ 1031--1084. Supplied report link: \url{https://doi.org/10.4171/owr/2023/19}. The record supplies no proposer or speaker name. Its verification fields label the statement \texttt{corrected\_verified} and say it was reconstructed and checked against the official report; those are supplied provenance claims, not verification performed for this draft.
\item Larsen--Tiep, \emph{Uniform character bounds for finite classical groups}, supplied literature link: \url{https://annals.math.princeton.edu/2024/200-1/p01}.
\item Supplied literature source: \emph{Oberwolfach Report 19/2023}, \url{https://doi.org/10.4171/owr/2023/19}.
\item Literature assessment dated exactly \texttt{2026-08-22}: the supplied curation labels the problem open and states that optimal uniform bounds for all finite almost quasisimple groups remain unresolved, while modern uniform bounds address important families and applications without providing the full explicit optimum sought here. This is supplied dated curation, not an independently verified literature assessment. The background explicitly requests independent verification of time-sensitive status and cited claims.
\item The excerpt attributes the reported \(79/80\) bound to Gluck [Gl], Gluck--Magaard [5], and Guralnick with the unnamed speaker [9], and an auxiliary classical-group estimate to Larsen, Shalev, and the unnamed speaker [11]. Complete bibliographic entries for these labels are absent from the supplied record.
\item The excerpt lists applications through Frobenius character formulas: the Ore conjecture [14] and Waring-type problems [11, 12]; random generation and representation varieties \(\operatorname{Hom}(\Gamma,G)\) for Fuchsian \(\Gamma\), with symmetric groups or finite-field groups arising from simple algebraic groups [17]; random walks on Cayley and McKay graphs [18, 19]; and Thompson's conjecture [1], stated there as the existence, for each finite nonabelian simple \(G\), of a conjugacy class \(C\) with \(G=C^2\). These are supplied contextual references, not assertions here about their current status. Their bibliography entries and some finite-field notation are missing from the transcription.
\end{itemize}
\end{document}
